\documentclass[10pt,a4paper]{amsart}
\usepackage{amsmath,amssymb,amsthm,graphicx}
% Retained preamble for the Stacks Project chapter {CH} (shared project preamble).
% The chapter file pinned-original-source.tex only carries \input{{preamble}} on its first line,
% so no preamble text lives inside the chapter itself.  The Stacks Project uses one single shared
% preamble for every chapter: the snapshot stages it three times (sources/stacks-chow/preamble.tex,
% sources/stacks-etale-cohomology/preamble.tex and focused-seeds/AGP/sources/stacks-snapshot-20261011/
% preamble.tex) and those copies are identical in size.  The block below is copied mechanically,
% line for line, from sources/stacks-chow/preamble.tex lines 14-18 and 154-217: the Xy-pic setup and
% the native theorem declarations together with the mathematical macros.  Nothing was added, removed
% or redefined; the wrapper class and the amsmath/amssymb/amsthm/graphicx base packages are supplied
% by the assembler.
% For commutative diagrams we use Xy-pic
% Chapter: stacks-spaces-more-morphisms.tex -- More on Morphisms of Spaces.
\usepackage[all]{xy}

% We use 2cell for 2-commutative diagrams.
\xyoption{2cell}

% Theorem environments.
%
\theoremstyle{plain}
\newtheorem{theorem}[subsection]{Theorem}
\newtheorem{proposition}[subsection]{Proposition}
\newtheorem{lemma}[subsection]{Lemma}

\theoremstyle{definition}
\newtheorem{definition}[subsection]{Definition}
\newtheorem{example}[subsection]{Example}
\newtheorem{exercise}[subsection]{Exercise}
\newtheorem{situation}[subsection]{Situation}

\theoremstyle{remark}
\newtheorem{remark}[subsection]{Remark}
\newtheorem{remarks}[subsection]{Remarks}

\numberwithin{equation}{subsection}

% Macros
%
\def\lim{\mathop{\mathrm{lim}}\nolimits}
\def\colim{\mathop{\mathrm{colim}}\nolimits}
\def\Spec{\mathop{\mathrm{Spec}}}
\def\Hom{\mathop{\mathrm{Hom}}\nolimits}
\def\Ext{\mathop{\mathrm{Ext}}\nolimits}
\def\SheafHom{\mathop{\mathcal{H}\!\mathit{om}}\nolimits}
\def\SheafExt{\mathop{\mathcal{E}\!\mathit{xt}}\nolimits}
\def\Sch{\mathit{Sch}}
\def\Mor{\mathop{\mathrm{Mor}}\nolimits}
\def\Ob{\mathop{\mathrm{Ob}}\nolimits}
\def\Sh{\mathop{\mathit{Sh}}\nolimits}
\def\NL{\mathop{N\!L}\nolimits}
\def\CH{\mathop{\mathrm{CH}}\nolimits}
\def\proetale{{pro\text{-}\acute{e}tale}}
\def\etale{{\acute{e}tale}}
\def\QCoh{\mathit{QCoh}}
\def\Ker{\mathop{\mathrm{Ker}}}
\def\Im{\mathop{\mathrm{Im}}}
\def\Coker{\mathop{\mathrm{Coker}}}
\def\Coim{\mathop{\mathrm{Coim}}}

% Boxtimes
%
\DeclareMathSymbol{\boxtimes}{\mathbin}{AMSa}{"02}

%
% Macros for moduli stacks/spaces
%
\def\QCohstack{\mathcal{QC}\!\mathit{oh}}
\def\Cohstack{\mathcal{C}\!\mathit{oh}}
\def\Spacesstack{\mathcal{S}\!\mathit{paces}}
\def\Quotfunctor{\mathrm{Quot}}
\def\Hilbfunctor{\mathrm{Hilb}}
\def\Curvesstack{\mathcal{C}\!\mathit{urves}}
\def\Polarizedstack{\mathcal{P}\!\mathit{olarized}}
\def\Complexesstack{\mathcal{C}\!\mathit{omplexes}}
% \Pic is the operator that assigns to X its picard group, usage \Pic(X)
% \Picardstack_{X/B} denotes the Picard stack of X over B
% \Picardfunctor_{X/B} denotes the Picard functor of X over B
\def\Pic{\mathop{\mathrm{Pic}}\nolimits}
\def\Picardstack{\mathcal{P}\!\mathit{ic}}
\def\Picardfunctor{\mathrm{Pic}}
\def\Deformationcategory{\mathcal{D}\!\mathit{ef}}
\usepackage{cleveref}
\crefname{originalreference}{original reference}{original references}
\Crefname{originalreference}{Original reference}{Original references}
\makeatletter
\newlabel{perfect-definition-relatively-perfect}{{0DI0}{source}{}{sourcepointer-1}{}}
\newlabel{perfect-definition-relatively-perfect@cref}{{[originalreference][0][]0DI0}{[1][0][]source}{}{}{}}
\makeatother
\crefname{originalreference}{original reference}{original references}
\Crefname{originalreference}{Original reference}{Original references}
\makeatletter
\newlabel{perfect-lemma-affine-locally-rel-perfect}{{0DI2}{source}{}{sourcepointer-2}{}}
\newlabel{perfect-lemma-affine-locally-rel-perfect@cref}{{[originalreference][0][]0DI2}{[1][0][]source}{}{}{}}
\makeatother
\crefname{originalreference}{original reference}{original references}
\Crefname{originalreference}{Original reference}{Original references}
\makeatletter
\newlabel{spaces-perfect-lemma-derived-quasi-coherent-small-etale-site}{{071Q}{source}{}{sourcepointer-3}{}}
\newlabel{spaces-perfect-lemma-derived-quasi-coherent-small-etale-site@cref}{{[originalreference][0][]071Q}{[1][0][]source}{}{}{}}
\makeatother
\crefname{originalreference}{original reference}{original references}
\Crefname{originalreference}{Original reference}{Original references}
\makeatletter
\newlabel{spaces-perfect-lemma-descend-pseudo-coherent}{{08HE}{source}{}{sourcepointer-4}{}}
\newlabel{spaces-perfect-lemma-descend-pseudo-coherent@cref}{{[originalreference][0][]08HE}{[1][0][]source}{}{}{}}
\makeatother
\crefname{originalreference}{original reference}{original references}
\Crefname{originalreference}{Original reference}{Original references}
\makeatletter
\newlabel{spaces-perfect-lemma-tor-dimension-rel}{{0DK7}{source}{}{sourcepointer-5}{}}
\newlabel{spaces-perfect-lemma-tor-dimension-rel@cref}{{[originalreference][0][]0DK7}{[1][0][]source}{}{}{}}
\makeatother
\crefname{originalreference}{original reference}{original references}
\Crefname{originalreference}{Original reference}{Original references}
\makeatletter
\newlabel{spaces-perfect-remark-match-total-direct-images}{{08GH}{source}{}{sourcepointer-6}{}}
\newlabel{spaces-perfect-remark-match-total-direct-images@cref}{{[originalreference][0][]08GH}{[1][0][]source}{}{}{}}
\makeatother
\title{Etale-local characterization of relatively perfect complexes on algebraic spaces}
\author{The Stacks Project Authors}
\date{stacks-project:spaces-more-morphisms.tex:local-snapshot-20261011; GNU FDL 1.2 or later}
\begin{document}
\maketitle
\noindent\textit{Editorial note.} Counted span is the author's own proof environment, source lines 12338-12364, immediately following the statement at 12312-12336.\par\medskip
\noindent\textit{Editorial note.} Clause (4)'s equivalence is imported from the schemes chapter by citation (proof lines 12365-12367); the reviewer should read that import against the bridge lemmas recorded in standard\_prerequisites.\par\medskip
\noindent\textit{Editorial note.} Four cross-chapter pointer families (Derived Categories of Spaces x3, Derived Categories of Schemes x1) are recorded in assembly.external\_references.\par\medskip
\noindent\textit{Allowed standard prerequisites (editorial).}\begin{itemize}
\item Small etale site versus scheme underlying an algebraic space. Derived Categories of Spaces, Lemma spaces-perfect-lemma-derived-quasi-coherent-small-etale-site, cited at proof lines 12347-12349. Identifies E|\_U in D\_QCoh(O\_U) with an object E\_0 of D\_QCoh(O\_\{U\_0\}) on the underlying schemes, the bridge on which all four clauses are compared.
\item Pseudo-coherence descends across the bridge. Derived Categories of Spaces, Lemma spaces-perfect-lemma-descend-pseudo-coherent, cited at proof lines 12350-12351. E|\_U is pseudo-coherent if and only if E\_0 is, for clauses (2)/(3) to speak the scheme language.
\item Relative tor dimension across the bridge. Derived Categories of Spaces, Lemma spaces-perfect-lemma-tor-dimension-rel, cited at proof lines 12352-12356. Finite tor dimension over g\textasciicircum{}\{-1\}O\_V transfers to finite tor dimension of E\_0 over g\_0\textasciicircum{}\{-1\}O\_\{V\_0\}.
\item Affine-local characterization on schemes. Derived Categories of Schemes, Lemma perfect-lemma-affine-locally-rel-perfect, cited at proof lines 12365-12367. The equivalence of clause (4) with clauses (2) and (3) is exactly the scheme-side statement, imported verbatim.
\item Etale locality of pseudo-coherence and of finite tor dimension. standard facts invoked at proof lines 12359-12361 (pseudo-coherence etale local; locally finite tor dimension etale local). Makes Y-perfectness testable etale locally, which powers (1) => (2) and (3) => (1).
\end{itemize}\medskip
\ifdefined\hypertarget\hypertarget{sourcepointer-1}{}\fi
\noindent\textit{Source reference.} Original reference 0DI0: focused-seeds/geometry/S9-spaces-affine-locally-rel-perfect/ref-tags.txt, source lines 5--5. Navigation pointer only: the counted proof cites this statement by \textbackslash{}ref as a standing input; the deduction the counted proof has to make is carried out inside the counted span, and no part of that deduction is delegated to this pointer.\par\medskip
\ifdefined\hypertarget\hypertarget{sourcepointer-2}{}\fi
\noindent\textit{Source reference.} Original reference 0DI2: focused-seeds/geometry/S9-spaces-affine-locally-rel-perfect/ref-tags.txt, source lines 6--6. Navigation pointer only: the counted proof cites this statement by \textbackslash{}ref as a standing input; the deduction the counted proof has to make is carried out inside the counted span, and no part of that deduction is delegated to this pointer.\par\medskip
\ifdefined\hypertarget\hypertarget{sourcepointer-3}{}\fi
\noindent\textit{Source reference.} Original reference 071Q: focused-seeds/geometry/S9-spaces-affine-locally-rel-perfect/ref-tags.txt, source lines 7--7. Navigation pointer only: the counted proof cites this statement by \textbackslash{}ref as a standing input; the deduction the counted proof has to make is carried out inside the counted span, and no part of that deduction is delegated to this pointer.\par\medskip
\ifdefined\hypertarget\hypertarget{sourcepointer-4}{}\fi
\noindent\textit{Source reference.} Original reference 08HE: focused-seeds/geometry/S9-spaces-affine-locally-rel-perfect/ref-tags.txt, source lines 8--8. Navigation pointer only: the counted proof cites this statement by \textbackslash{}ref as a standing input; the deduction the counted proof has to make is carried out inside the counted span, and no part of that deduction is delegated to this pointer.\par\medskip
\ifdefined\hypertarget\hypertarget{sourcepointer-5}{}\fi
\noindent\textit{Source reference.} Original reference 0DK7: focused-seeds/geometry/S9-spaces-affine-locally-rel-perfect/ref-tags.txt, source lines 9--9. Navigation pointer only: the counted proof cites this statement by \textbackslash{}ref as a standing input; the deduction the counted proof has to make is carried out inside the counted span, and no part of that deduction is delegated to this pointer.\par\medskip
\ifdefined\hypertarget\hypertarget{sourcepointer-6}{}\fi
\noindent\textit{Source reference.} Original reference 08GH: focused-seeds/geometry/S9-spaces-affine-locally-rel-perfect/ref-tags.txt, source lines 10--10. Navigation pointer only: the counted proof cites this statement by \textbackslash{}ref as a standing input; the deduction the counted proof has to make is carried out inside the counted span, and no part of that deduction is delegated to this pointer.\par\medskip
\begin{lemma}
\label{lemma-affine-locally-rel-perfect}
Let $S$ be a scheme. Let $f : X \to Y$ be a morphism of algebraic spaces
over $S$ which is flat and locally of finite presentation.
Let $E \in D_\QCoh(\mathcal{O}_X)$. The following are equivalent:
\begin{enumerate}
\item $E$ is $Y$-perfect,
\item for every commutative diagram
$$
\xymatrix{
U \ar[d] \ar[r]_g & V \ar[d] \\
X \ar[r]^f & Y
}
$$
where $U$, $V$ are schemes and the vertical arrows are \'etale, the complex
$E|_U$ is $V$-perfect in the sense of Derived Categories of Schemes,
Definition \ref{perfect-definition-relatively-perfect},
\item for some commutative diagram as in (2) with $U \to X$
surjective, the complex $E|_U$ is $V$-perfect in the sense of
Derived Categories of Schemes,
Definition \ref{perfect-definition-relatively-perfect},
\item for every commutative diagram as in (2) with $U$ and $V$
affine the complex $R\Gamma(U, E)$ is $\mathcal{O}_Y(V)$-perfect.
\end{enumerate}
\end{lemma}

\begin{proof}
To make sense of parts (2), (3), (4) of the lemma, observe that
the object $E|_U$ of $D_\QCoh(\mathcal{O}_U)$ corresponds to
an object $E_0$ of $D_\QCoh(\mathcal{O}_{U_0})$ where $U_0$
denotes the scheme underlying $U$, see Derived Categories of Spaces,
Lemma \ref{spaces-perfect-lemma-derived-quasi-coherent-small-etale-site}.
Moreover, in this case $E_0$ is pseudo-coherent if and only if
$E|_U$ is pseudo-coherent, see Derived Categories of Spaces,
Lemma \ref{spaces-perfect-lemma-descend-pseudo-coherent}.
Also, $E|_U$ locally has finite tor dimension over
$f^{-1}\mathcal{O}_Y|_U = g^{-1}\mathcal{O}_V$ if and only if
$E_0$ locally has finite tor dimension over $g_0^{-1}\mathcal{O}_{V_0}$
by Derived Categories of Spaces, Lemma
\ref{spaces-perfect-lemma-tor-dimension-rel}.
Here $g_0 : U_0 \to V_0$ is the morphism of schemes representing
$g : U \to V$ (notation as in Derived Categories of Spaces,
Remark \ref{spaces-perfect-remark-match-total-direct-images}).
Finally, observe that ``being pseudo-coherent'' is \'etale local and
of course ``having locally finite tor dimension'' is \'etale local.
Thus we see that it suffices to check $Y$-perfectness \'etale locally
and by the above discussion we see that (1) implies (2) and
(3) implies (1). Since part (4) is equivalent
to (2) and (3) by
Derived Categories of Schemes, Lemma
\ref{perfect-lemma-affine-locally-rel-perfect}
the proof is complete.
\end{proof}

\end{document}
